% =====================================================================
% paper/sections/critic_degeneracy_hypothesis.tex
%
% PROVENANCE. Berkeley B-SYNTH row 12 (iter144, 2026-07-04):
% critic-degeneracy hypothesis (CDH) empirical test on the same-stack
% PPO/GRPO benchmark. Driver: scripts/berkeley/critic_degeneracy_hypothesis.py
% Inputs:  experiments/results/samestack_ppo_grpo.json
%          experiments/results/berkeley/eureka_rqs_per_anchor.tsv
% Outputs: experiments/results/berkeley/cdh_{gradnorm_stats,reward_window,
%          paired_test,gradnorm_vs_reward,rqs_collapse}.tsv
%          experiments/results/berkeley/cdh_summary.json
%
% Frontier synthesis (FRONTIER_INSIGHTS.md Round 1, ChatGPT Pro
% Extended + Gemini Deep Think) proposes the Critic-Degeneracy
% Hypothesis: PPO's value head V_phi(x_{1:t}) is mathematically
% degenerate for sparse, terminal-reward CoT -- it collapses to a
% static prompt-difficulty regressor V_phi(x_{1:t}) approx E[R|x_prompt],
% what GRPO computes statelessly via the group mean. Under CDH the
% critic is a NET NOISE source, not a variance reducer. We test the
% OBSERVABLE consequences on the same-stack benchmark.
% =====================================================================

\subsection{Frontier-Model Cross-Examination: The Critic-Degeneracy Hypothesis}
\label{sec:frontier-critic-degeneracy}

To stress-test the PPO$\leftrightarrow$GRPO equivalence from
Section~\ref{sec:group-size} we ran an empirical follow-on suggested
by an external frontier-model cross-examination (FRONTIER\_INSIGHTS.md
Round~1, ChatGPT Pro Extended + Gemini Deep Think). The proposal is
the \textbf{Critic-Degeneracy Hypothesis} (CDH): for sparse,
terminal-reward chain-of-thought, PPO's value head
$V_\phi(x_{1:t})$ is mathematically degenerate---it collapses to a
static prompt-difficulty regressor
$V_\phi(x_{1:t}) \approx E[R \mid x_\text{prompt}]$, i.e.\ the very
scalar that GRPO computes \emph{statelessly} via the group mean. If
CDH holds, the PPO critic is a noise \emph{amplifier} (its parametric
fit injects variance around the prompt-mean) rather than the variance
\emph{reducer} the literature typically assumes.

\paragraph{Scope.}
This is a small, short-horizon test: 40 optimizer steps on a 0.5B
model on near-saturated arithmetic, with a cold-start critic (a
randomly initialised single linear layer read off the prompt-end
hidden state). It cannot distinguish asymptotic critic degeneracy
from a merely unconverged critic, and related work cuts both ways
outside this regime: \citet{kazemnejad2025vineppo} find
Monte-Carlo value estimates outperform a learned PPO critic on
reasoning tasks, while \citet{bytedance2025vapo} show that a
pretrained, GAE-calibrated value model can outperform critic-free
baselines on hard reasoning. CDH is tested here only as scoped to
this cold-start, short-horizon setup.

\paragraph{Five hypotheses.}
We test the observable consequences of CDH on the same-stack
benchmark (5 seeds $\times$ 2 algos, Qwen2.5-0.5B / arithmetic, 40
steps, $k{=}2$ inner epochs). The ``128'' in the run logs is
generations \emph{per optimizer step}, not a group size: GRPO
draws 16 prompts $\times$ group size 8, PPO draws 128 distinct
prompts $\times\,1$:
\begin{enumerate}\itemsep2pt
\item \textbf{H1 -- CV(grad\_norm).} If the critic is a control
  variate, CV(grad\_norm) should be \emph{lower} for PPO than GRPO.
  We observe GRPO CV$=$1.347 vs PPO CV$=$1.433: PPO is $6\%$ more
  variable in gradient norm---the opposite of the variance-reduction
  prediction, though at $n{=}5$ seeds this $6\%$ CV difference
  carries no confidence interval or significance test and is
  descriptive only. The raw gradient norms differ sharply in level:
  PPO $E\|g\|\,{=}\,96.79$ vs GRPO $E\|g\|\,{=}\,0.62$, a
  $156\times$ ratio on the same stack. The logged
  \texttt{grad\_norm} is the pre-clipping total norm returned by
  \texttt{clip\_grad\_norm\_} (clip value $1.0$) on the final inner
  epoch; for PPO the parameter set includes the value head and the
  loss includes a $0.5\times$ value-MSE term, so the level gap
  partly reflects the critic's own loss rather than actor-gradient
  scale alone---an uncontrolled confound for the level comparison
  (though not for the within-algorithm CV).
\item \textbf{H2 -- rolling-variance of mean\_reward (window 10).}
  If the critic smooths the trajectory, PPO should have lower
  within-run variance. We observe GRPO $0.0049$ vs PPO $0.0085$:
  PPO's per-step reward is $73\%$ noisier than GRPO's---directionally
  the opposite of the critic smoothing the trajectory.
\item \textbf{H3 -- paired last10\_avg (5 seeds, matched-stack).}
  CDH predicts near-equivalence (critic $\approx$ group mean), but
  this test cannot certify it. The paired differences give mean
  $\Delta_{\text{GRPO-PPO}}=0.0608$, $sd{=}0.0958$, $n{=}5$, hence
  $\mathrm{SE}=0.0958/\sqrt{5}=0.0428$ and a $95\%$ confidence
  interval of approximately $[-0.058,\,0.180]$ (paired $t$,
  $\nu{=}4$; two-sided $p{=}0.156$). The result is inconclusive:
  the interval contains both $0$ and effects beyond $\pm0.1$, and
  the $6$-point point estimate itself exceeds an
  $\epsilon{=}0.02$ TOST equivalence margin, so no equivalence
  claim is licensed by failing to reject the null.
\item \textbf{H4 -- $\mathrm{corr}(\|g\|,\,R_\text{batch})$.} If the
  critic reduced variance, PPO's gradient norm should track batch
  reward at least as tightly as GRPO's; the CDH/noise reading
  predicts the opposite---a \emph{looser} coupling (smaller $|r|$)
  for PPO. We observe GRPO $r{=}{-}0.553$ vs PPO
  $r{=}{-}0.445$ (per-seed mean): PPO's $|r|$ is $19.5\%$ smaller
  than GRPO's (base: GRPO), equivalently GRPO's $|r|$ is $24\%$
  larger than PPO's (base: PPO). This is directionally consistent
  with CDH, but it is also mechanically confounded: GRPO's centred
  gradient vanishes as a batch saturates ($\|g\|\to0$ as
  $R\to1$ through the zero-variance floor), which induces a
  negative $\|g\|$--$R$ coupling for GRPO by construction.
\item \textbf{H5 -- regressor-collapse fingerprint.} Under CDH the
  PPO critic is trying to learn $E[R \mid x_\text{prompt}]$, which
  RQS (Reward-Design Quality Score: a one-dimensional,
  non-parametric per-prompt difficulty score in $[0,1]$, released
  per prompt anchor in
  \texttt{experiments/results/berkeley/eureka\_rqs\_per\_anchor.tsv})
  approximates without parameters. If CDH holds, an OLS regression
  of the \emph{empirical prompt-mean reward} $r_\text{mean}$---the
  Monte-Carlo mean reward per anchor prompt, \emph{not} measured
  token-level critic predictions $V_\phi(x_{1:t})$, which this test
  does not observe---on RQS should explain substantial variance.
  Across the 7 prompt anchors in the released TSV with
  $0 < \mathrm{RQS} < 1$:
  \begin{equation*}
    \hat r_\text{mean} = -0.082 + 1.282 \cdot \mathrm{RQS},\quad
    R^2 = 0.490,\quad n = 7.
  \end{equation*}
  A 1-dimensional static regressor explains \emph{half} of
  $r_\text{mean}$ variance. A PPO value head (typically
  $L \cdot d_h^2$ parameters) would be fitting residual variation
  on top of this 1-parameter axis.
\end{enumerate}

\paragraph{Verdict.}
H1, H2, and H4 point against the variance-reduction reading of
PPO's critic in this setup (descriptively, at $n{=}5$, and with
the confounds noted under H1 and H4), and H5 shows a
one-dimensional static regressor captures half the prompt-mean
reward variance. H3 is inconclusive, and it must not be conflated
with the held-out comparison: the last-10-step paired difference
($\Delta{=}0.0608$, $p{=}0.156$) and the held-out paired
difference ($\Delta{=}{-}0.002$, $p{=}0.62$) are different
quantities, and \emph{neither} certifies equivalence---both are
under-powered at $n{=}5$. Taken together, the results are
\emph{consistent with} the Critic-Degeneracy Hypothesis in this
small cold-start setup: the critic behaves like a noisy
re-parameterisation of the prompt mean that GRPO computes
statelessly, and which RQS approximates one-dimensionally with
$R^2{=}0.49$. They do not establish CDH as a general property of
outcome-reward RL.

\paragraph{Implications for Pillar~1 and Pillar~3.}
Pillar~1's same-stack held-out comparison (paired
$\Delta_{\text{GRPO-PPO}} = -0.002$, $p{=}0.62$) is
\emph{consistent with} a CDH reading---a critic fitting the
prompt mean would produce near-identical updates to GRPO's
stateless baseline---without certifying equivalence at $n{=}5$.
Pillar~3's G-sweep result (G${=}2$ retention of G${=}16$ is
$100.3\%$ on near-ceiling arithmetic) admits the same reading:
in this saturated regime the prompt mean is the dominant signal
axis, so rollouts beyond what already pins it contribute little.
The static RQS regressor (H5) is a one-dimensional summary of
that axis; whether the PPO critic and the GRPO group mean remain
interchangeable outside this small cold-start setup is exactly
what CDH, as scoped here, does not decide.